\section*{الفصل \ref{ch:b3:group-theory-applied} --- تطبيقات نظرية الزمر}

\begin{solution}{exo:b3:group-theory-applied:1}
$n_{A_1} = \frac14(5 - 1 + 3 + 1) = 2$، $n_{A_2} = \frac14(5 - 1 - 3 - 1) = 0$،
$n_{B_1} = \frac14(5 + 1 + 3 - 1) = 2$، $n_{B_2} = \frac14(5 + 1 - 3 + 1) = 1$:
$\Gamma = 2A_1 + 2B_1 + B_2$، وبعده 5.
\end{solution}

\begin{solution}{exo:b3:group-theory-applied:2}
عدّ الذرات الباقية في مواضعها هو نفسه كما في الماء: $\Gamma_{\mathrm{vib}} = 2A_1 + B_1$. النمطان $A_1$ ($z$) 
و$B_1$ ($x$) نشطان في تحت الأحمر، وكلاهما \omterm{def:b3:group-theory-applied:activity}{نشط في رامان} أيضًا ($x^2$، $xz$): ثلاث
حزم في كل طيف.
\end{solution}

\begin{solution}{exo:b3:group-theory-applied:3}
$B_1\otimes B_2$: \omterm{def:b3:point-groups:representation}{المميِّزات} $(1, 1, -1, -1)$ $= A_2$. $A_2\otimes B_1$: $(1, -1, -1,
1) = B_2$. وفي $C_{3v}$، للجداء $E\otimes E$ \omterm{def:b3:point-groups:representation}{المميِّزات} $(4, 1, 0)$؛ والاختزال: $A_1 + A_2 +
E$.
\end{solution}

\begin{solution}{exo:b3:group-theory-applied:4}
تحت الأحمر: النمطان $t_2$، عند 3019 و\qty{1306}{cm^{-1}}. رامان: الأربعة كلها ($a_1$
و$e$ و$t_2$ تتحول كلها كما تتحول دوال تربيعية). ولا \omterm{def:b3:point-groups:inversion}{مركز انقلاب}، فالتطابقات
بين الطيفين مسموحة.
\end{solution}

\begin{solution}{exo:b3:group-theory-applied:5}
عدد الذرات الباقية في مواضعها هو $(4, 1, 2, 4, 1, 2)$، وبضربه في $\chi_{xyz} = (3, 0, -1, 1, -2, 1)$ نحصل على
$\chi_{3N} = (12, 0, -2, 4, -2, 2)$، الذي يُختزل إلى $A_1' + A_2' + 3E' + 2A_2'' +
E''$. وبطرح $E' + A_2''$ (الانسحابات) و$A_2' + E''$ (الدورانات): $A_1' + 2E'
+ A_2''$. تحت الأحمر: $2E' + A_2''$، ثلاث حزم؛ ورامان: $A_1' + 2E'$، ثلاث حزم. والنمط $A_2''$
(الانحناء خارج المستوي) في تحت الأحمر وحده، و$A_1'$ (التمدد المتناظر) في رامان وحده، 
والنمطان $E'$ يظهران في كليهما.
\end{solution}

\begin{solution}{exo:b3:group-theory-applied:6}
تحت الأحمر: النمطان $T_{1u}$، حزمتان. رامان: $A_{1g}$ و$E_g$ و$T_{2g}$، ثلاث حزم.
والنمط $T_{2u}$ صامت. ولجزيء \ce{SF6} \omterm{def:b3:point-groups:inversion}{مركز انقلاب}: الاستبعاد المتبادل.
\end{solution}

\begin{solution}{exo:b3:group-theory-applied:7}
تحت $E$ و$C_3$ و$C_3^2$ تنتقل الدالة $h_2$ إلى $h_2$ و$h_3$ و$h_1$ (\omterm{def:b3:point-groups:representation}{المميِّزات} 2
و$-1$ و$-1$)؛ والانعكاسات تساهم بصفر. $\hat P_Eh_2 \propto 2h_2 - h_3 - h_1$.
وتداخله مع $2h_1 - h_2 - h_3$ (بإهمال التداخلات الذرية) $-3$، 
ولهذه الأخيرة المعيار $\sqrt6$؛ وبطرح الإسقاط، $(2h_2 - h_1 - h_3) +
\frac12(2h_1 - h_2 - h_3) = \frac32(h_2 - h_3)$.
\end{solution}

\begin{solution}{exo:b3:group-theory-applied:8}
fac ($C_{3v}$، \omterm{def:b3:point-groups:representation}{المميِّزات} $3, 0, 1$): $A_1 + E$، حزمتان في تحت الأحمر. mer ($C_{2v}$،
\omterm{def:b3:point-groups:representation}{المميِّزات} $3, 1, 3, 1$): $2A_1 + B_1$، ثلاث حزم. \ce{M(CO)6} ($O_h$): $A_{1g} + E_g
+ T_{1u}$، حزمة واحدة في تحت الأحمر ($T_{1u}$).
\end{solution}

\begin{solution}{exo:b3:group-theory-applied:9}
\omterm{def:b3:point-groups:representation}{المميِّزات} $6, 0, 0, 2, 2, 0, 0, 0, 4, 2$: لا يُبقي ربائطَ في مواضعها إلا $E$، والمحاور $C_4$ و$C_2$ الواقعة على امتداد
المحاور الإحداثية (ربيطتان باقيتان في موضعيهما)، والمستويات الثلاثة $\sigma_h$ (أربع) والمستويات الستة $\sigma_d$
(اثنتان). ويعطي الاختزال $A_{1g} + E_g + T_{1u}$. فتتحد مدارات الفلز $s$
($a_{1g}$) و$d_{z^2}$ و$d_{x^2-y^2}$ ($e_g$) و$p_x$ و$p_y$ و$p_z$ ($t_{1u}$)؛
وتبقى $t_{2g}$ غير رابطة.
\end{solution}

\begin{solution}{exo:b3:group-theory-applied:10}
$n_i = \frac1h(1\cdot h\cdot\chi_i(E)) = d_i$. وبعد التمثيل المنتظم
هو $h$، وهو أيضًا $\sum_in_id_i = \sum_id_i^2$.
\end{solution}

\begin{solution}{exo:b3:group-theory-applied:11}
في $D_{2h}$ (والجزيء على امتداد $z$)، الذرات الباقية في مواضعها $(4, 4, 0, 0, 0, 0, 4, 4)$، و$\chi_{3N}
= (12, -4, 0, 0, 0, 0, 4, 4)$، الذي يُختزل إلى $2A_g + 2B_{2g} + 2B_{3g} + 2B_{1u} + 2B_{2u}
+ 2B_{3u}$. وبطرح الانسحابات $B_{1u} + B_{2u} + B_{3u}$ والدورانين
$B_{2g} + B_{3g}$: $2A_g + B_{2g} + B_{3g} + B_{1u} + B_{2u} + B_{3u}$، أي سبعة أنماط.
وفي $D_{\infty h}$: $2\Sigma_g^+ + \Pi_g + \Sigma_u^+ + \Pi_u$. تحت الأحمر: $\Sigma_u^+$ و$\Pi_u$
(حزمتان)؛ ورامان: $2\Sigma_g^+$ و$\Pi_g$ (ثلاث حزم)؛ ولا تطابق.
\end{solution}

\begin{solution}{exo:b3:group-theory-applied:12}
$1b_2 \to 4a_1$: $B_2$، $y$، مسموح، عمودي على المستوي. $3a_1 \to 4a_1$:
$A_1$، $z$، مسموح، على امتداد المحور $C_2$. $1b_2 \to 2b_1$: $A_2$، ممنوع.
$1b_1 \to 4a_1$: $B_1$، $x$، مسموح، في المستوي، عموديًا على المحور.
\end{solution}

\begin{solution}{pb:b3:group-theory-applied:1}
\textbf{1.} cis: مجموعتا CO على رأسين متجاورين؛ trans: متقابلين؛ fac: ثلاث
مجموعات CO على وجه واحد؛ mer: ثلاث مجموعات CO على خط طول (اثنتان في وضع trans، وواحدة بينهما).
\textbf{2.} $C_{2v}$.
\textbf{3.} $D_{4h}$.
\textbf{4.} $C_{3v}$.
\textbf{5.} $C_{2v}$.
\textbf{6.} trans-\ce{ML4(CO)2} وحده.
\textbf{7.} العملية التي تنقل CO إلى أخرى تضع متجهته خارج
القطر (مساهمة 0)؛ والتي تُبقي CO في موضعها تنقل متجهة رابطتها
إلى نفسها (+1)؛ ولا عملية تعكس متجهة C--O.
\textbf{8.} \omterm{def:b3:point-groups:representation}{المميِّزات} $(2, 0, 2, 0)$، مع $\sigma_v(xz)$ مستوي مجموعتي CO كلتيهما: $A_1 +
B_1$.
\textbf{9.} \omterm{def:b3:point-groups:representation}{المميِّزات} 2 تحت $E$ و$C_4$ و$C_2$ و$\sigma_v$ و$\sigma_d$، و0 في غيرها:
$A_{1g} + A_{2u}$.
\textbf{10.} $(3, 0, 1)$: $A_1 + E$.
\textbf{11.} $(3, 1, 3, 1)$: $2A_1 + B_1$.
\textbf{12.} 2، 2، 3، 3: بعد واحد لكل CO.
\textbf{13.} cis 2، trans 1، fac 2، mer 3.
\textbf{14.} cis 2، trans 1، fac 2، mer 3.
\textbf{15.} trans: حزمته في تحت الأحمر ($A_{2u}$) وحزمته في رامان ($A_{1g}$) نمطان
مختلفان.
\textbf{16.} يغيّر $A_{2u}$ إشارته تحت $i$: تطول إحدى مجموعتي CO بينما تقصر
الأخرى، بطورين متعاكسين.
\textbf{17.} $E$ زوج منحل: نمطان لهما التردد نفسه يعطيان
حزمة واحدة، تضاف إليها حزمة $A_1$.
\textbf{18.} متماكب ثنائي الاستبدال في وضع trans، متناظر مركزيًا.
\textbf{19.} ثلاث حزم في تحت الأحمر: mer (فالمتماكب fac يعطي حزمتين).
\textbf{20.} نعم: للمتماكب mer ثلاثة أنماط CO نشطة في رامان عند الترددات نفسها،
لأنه غير متناظر مركزيًا؛ أما fac فيُظهر حزمتين.
\textbf{21.} $A_1$: التمدد المتطاور متناظر كليًا.
\textbf{22.} خليط من fac وmer يُظهر حتى خمس حزم في تحت الأحمر، مع حزم fac
في مواضع مختلفة؛ أما ثلاث حزم نظيفة فتوافق متماكبًا واحدًا. وطيف
الرنين المغناطيسي النووي \babelsublr{${}^{13}$C} يحسم الأمر: إشارتان للكربونيل بنسبة 2:1 للمتماكب
mer، وإشارة واحدة للمتماكب fac.
\textbf{23.} \textbf{للمتماكب mer ثلاثة تمددات CO نشطة في تحت الأحمر، وللمتماكب fac
تمددان}: فالناتج هو mer.
\end{solution}
